\section*{Capítulo \ref{ch:b1:optical-mineralogy} --- Mineralogía óptica}

\begin{solution}{exo:b1:optical-mineralogy:1}
Cuarzo: $\qty{30000}{nm}\times0.009 = \qty{270}{nm}$, un blanco de primer \omterm{def:b1:optical-mineralogy:interference-colour}{orden}. Calcita:
$\qty{30000}{nm}\times0.172 = \qty{5160}{nm}$, muy por encima del diagrama: un blanco pálido de \omterm{def:b1:optical-mineralogy:interference-colour}{orden} alto,
ligeramente rosado o verdoso.
\end{solution}

\begin{solution}{exo:b1:optical-mineralogy:2}
$I/I_{\max} = \sin^2 2\theta$: $0$ en \ang{0}; $\sin^2\ang{45} = 0.5$ en \ang{22.5}; $1$ en \ang{45}; $0.5$ en
\ang{67.5}; $0$ en \ang{90}.
\end{solution}

\begin{solution}{exo:b1:optical-mineralogy:3}
El \omterm{def:b1:rock-forming-minerals:groups}{granate} es cúbico, luego isótropo: la luz sale de él vibrando como entró, según el polarizador, y el analizador la
detiene; $\Delta = 0$ en cualquier posición. Un grano de cuarzo cortado perpendicularmente a su \omterm{def:b1:optical-mineralogy:optic-axis}{eje óptico} tiene también
$n_1 = n_2$ para la luz que viaja según ese eje, luego $\Delta = 0$: también permanece oscuro. En \omterm{def:b1:field-lab-methods-1:microscope}{luz polarizada plana} el
\omterm{def:b1:rock-forming-minerals:groups}{granate} tiene \omterm{def:b1:optical-mineralogy:relief}{relieve} muy alto ($1.83$ frente a $1.54$), el cuarzo casi ninguno.
\end{solution}

\begin{solution}{exo:b1:optical-mineralogy:4}
$|n - 1.54|$: almandino $0.29$, forsterita $0.11$, ortoclasa $0.02$, cuarzo $0.01$; \omterm{def:b1:optical-mineralogy:relief}{relieve} en ese \omterm{def:b1:optical-mineralogy:interference-colour}{orden}, de mayor a menor.
La ortoclasa tiene el índice menor: \omterm{def:b1:optical-mineralogy:relief}{relieve} negativo, y su \omterm{def:b1:optical-mineralogy:relief}{línea de Becke} sale del grano, hacia el medio de montaje, al
subir el foco.
\end{solution}

\begin{solution}{exo:b1:optical-mineralogy:5}
$\delta = 650/30000 = 0.022$. Un \omterm{def:b1:rock-forming-minerals:mineral}{mineral} verde \omterm{def:b1:optical-mineralogy:pleochroism}{pleocroico} con exfoliaciones a \ang{56} y \ang{124} y una \omterm{def:b1:optical-mineralogy:birefringence}{birrefringencia}
cercana a $0.020$: hornblenda.
\end{solution}

\begin{solution}{exo:b1:optical-mineralogy:6}
$\tan\ang{30} = 0.5774$. Cuarzo: $(0.013/1.54)\times0.5774 = \qty{0.00487}{rad}
= \ang{0.28}$. Forsterita: $(0.111/1.54)\times0.5774 = \qty{0.0416}{rad} =
\ang{2.38}$. Cociente $0.111/0.013 = 8.5$. Exactamente: $\sin r = 1.54\sin\ang{30}/1.651
= 0.4664$, $r = \ang{27.80}$, una desviación de \ang{2.20}: la fórmula de primer \omterm{def:b1:optical-mineralogy:interference-colour}{orden} sobrestima en un
\qty{8}{\percent} para una diferencia de índices tan grande.
\end{solution}

\begin{solution}{exo:b1:optical-mineralogy:7}
$270 + 540 = \qty{810}{nm}$, un verde de segundo \omterm{def:b1:optical-mineralogy:interference-colour}{orden}, cuando las direcciones lentas del grano y de la lámina son
paralelas (los \omterm{def:b1:optical-mineralogy:retardation}{retardos} se suman); $|270 - 540| =
\qty{270}{nm}$, de nuevo blanco de primer \omterm{def:b1:optical-mineralogy:interference-colour}{orden}, cuando la dirección
lenta del grano es paralela a la dirección rápida de la lámina.
\end{solution}

\begin{solution}{exo:b1:optical-mineralogy:8}
A espesor completo $\Delta = \qty{30000}{nm}\times0.052 = \qty{1560}{nm}$. Desde el borde ($\Delta = 0$) hacia dentro, el
\omterm{def:b1:optical-mineralogy:retardation}{retardo} pasa por 542 y \qty{1074}{nm}, pero no por \qty{1605}{nm}: dos franjas rojo violáceas.
\end{solution}

\begin{solution}{exo:b1:optical-mineralogy:9}
El \omterm{def:b1:optical-mineralogy:extinction}{ángulo de extinción} medido en un grano depende de cómo está cortado el grano: solo las secciones con una orientación
particular muestran el ángulo mayor, y las demás muestran menos, hasta cero. El valor mayor, unos \ang{37}, es el que hay
que comparar: coincide con el diópsido (unos \ang{38}), no con la hornblenda (como máximo \ang{19}).
\end{solution}

\begin{solution}{exo:b1:optical-mineralogy:10}
Con el analizador según $x$, las componentes según $x$ de las dos ondas son
$E_0\cos\theta\cdot\cos\theta$ y $-E_0\sin\theta\,e^{-i\varphi}\cdot(-\sin\theta)$, así que
$A = E_0(\cos^2\theta + \sin^2\theta\,e^{-i\varphi})$ y
\[
  |A|^2 = E_0^2\bigl(\cos^4\theta + \sin^4\theta + 2\sin^2\theta\cos^2\theta\cos\varphi\bigr)
  = E_0^2\bigl(1 - 2\sin^2\theta\cos^2\theta\,(1 - \cos\varphi)\bigr),
\]
usando $\cos^4\theta + \sin^4\theta = 1 - 2\sin^2\theta\cos^2\theta$. Como
$2\sin^2\theta\cos^2\theta = \frac12\sin^2 2\theta$ y $1 - \cos\varphi =
2\sin^2(\varphi/2)$, $I_\parallel = I_0(1 - \sin^2 2\theta\sin^2(\pi\Delta/\lambda))$.
A cada longitud de onda, $I_\parallel + I_\perp = I_0$: los colores entre polarizadores paralelos son los complementarios
de los que hay entre \omterm{def:b1:field-lab-methods-1:microscope}{nicoles cruzados}. Un grano de \qty{550}{nm}, magenta entre \omterm{def:b1:field-lab-methods-1:microscope}{nicoles cruzados}, es verde amarillento
pálido entre polarizadores paralelos.
\end{solution}

\begin{solution}{exo:b1:optical-mineralogy:11}
(a) $\Delta/\lambda$ va desde $5000/700 = 7.14$ hasta $5000/400 = 12.5$; pasa por los enteros 8 a 12: cinco ceros.
(b) Los ceros están en $\lambda_k = \Delta/k$; dos sucesivos están separados
$\Delta/k - \Delta/(k + 1) =
\Delta/(k(k + 1)) \approx \lambda^2/\Delta$, es decir, $550^2/5000 =
\qty{60.5}{nm}$ cerca de \qty{550}{nm}. (c) Cinco bandas oscuras y estrechas repartidas por el espectro quitan un poco de
cada color, y lo que queda pasa con un peso medio de la mitad: el ojo ve un tinte pálido, casi neutro.
\end{solution}

\begin{solution}{exo:b1:optical-mineralogy:12}
La media de $\sin^2 2\theta$ sobre una vuelta es $\frac1{2\pi}\int_0^{2\pi}\sin^2 2\theta\,
\mathrm d\theta = \frac12$, así que la intensidad media es $\frac12 I_0\sin^2(\pi\Delta/\lambda)$. Ella calcula
$\delta = \Delta/30000$ en lugar de $\Delta/25000$ (en nm): todos sus valores son $25/30 = 0.83$ de los verdaderos,
demasiado bajos en un \qty{17}{\percent}.
\end{solution}

\begin{solution}{pb:b1:optical-mineralogy:1}
\textbf{1.} B: los bordes oscuros indican \omterm{def:b1:optical-mineralogy:relief}{relieve} alto; los bordes apenas visibles de A indican \omterm{def:b1:optical-mineralogy:relief}{relieve} bajo.

\textbf{2.} A: $(1.544 + 1.553)/2 = 1.5485$, $|n - n_m| = 0.0085$. B:
$(1.664 + 1.672 + 1.694)/3 = 1.6767$, $|n - n_m| = 0.1367$.

\textbf{3.} Las desviaciones son proporcionales a $|n - n_m|$: $0.1367/0.0085 =
16.1$.

\textbf{4.} B tiene un índice mayor que el medio de montaje: \omterm{def:b1:optical-mineralogy:relief}{relieve} positivo.

\textbf{5.} \omterm{def:b1:optical-mineralogy:pleochroism}{Pleocroísmo}; biotita.

\textbf{6.} Dos exfoliaciones a \ang{56} y \ang{124}: un \omterm{def:b1:rock-forming-minerals:groups}{anfíbol}.

\textbf{7.} $\qty{30000}{nm}\times0.009 = \qty{270}{nm}$: blanco de primer \omterm{def:b1:optical-mineralogy:interference-colour}{orden}.

\textbf{8.} Cada grano se corta según su propia orientación y muestra un valor de $|n_1 - n_2|$ entre 0 y $\delta$;
solo los granos cortados paralelamente al \omterm{def:b1:optical-mineralogy:optic-axis}{eje óptico} muestran el propio $\delta$. El color más alto entre muchos granos
es el que hay que usar.

\textbf{9.} $d = 342/0.009 = \qty{38000}{nm} = \qty{38}{\micro m}$.

\textbf{10.} $(38 - 30)/30 = \qty{26.7}{\percent}$.

\textbf{11.} $\sin^2(\pi\times342/550) = 0.861$; a \qty{30}{\micro m},
$\sin^2(\pi\times270/550) = 0.999$.

\textbf{12.} $\qty{38000}{nm}\times0.008 = \qty{304}{nm}$ a
$\qty{38000}{nm}\times0.013 = \qty{494}{nm}$: de blanco a rojo anaranjado, todo por debajo del final del primer \omterm{def:b1:optical-mineralogy:interference-colour}{orden} a
\qty{542}{nm}.

\textbf{13.} Entre 1074 y \qty{1605}{nm}: tercer \omterm{def:b1:optical-mineralogy:interference-colour}{orden}.

\textbf{14.} $\delta = 1140/38000 = 0.030$.

\textbf{15.} $1140/30000 = 0.038$, lo que sugeriría un \omterm{def:b1:rock-forming-minerals:groups}{olivino}. Un \omterm{def:b1:rock-forming-minerals:groups}{olivino} magnésico y el cuarzo reaccionan para dar un
\omterm{def:b1:rock-forming-minerals:groups}{piroxeno}, \ce{Mg2SiO4 + SiO2 -> 2 MgSiO3}; una roca que ha cristalizado con cuarzo no debería conservar ambos.

\textbf{16.} $0.030$ es la \omterm{def:b1:optical-mineralogy:birefringence}{birrefringencia} del diópsido; su \omterm{def:b1:optical-mineralogy:relief}{relieve} alto, su falta de color y sus dos exfoliaciones cerca
de \ang{87} concuerdan: B es un clinopiroxeno, diópsido.

\textbf{17.} $\qty{30000}{nm}\times0.030 = \qty{900}{nm}$: amarillo de segundo \omterm{def:b1:optical-mineralogy:interference-colour}{orden}.

\textbf{18.} En \qty{8}{\micro m}; el cuarzo muestra entonces como mucho \qty{270}{nm}, blanco de primer \omterm{def:b1:optical-mineralogy:interference-colour}{orden}.

\textbf{19.} \omterm{def:b1:optical-mineralogy:twin}{Macla polisintética}: D es \omterm{def:b1:rock-forming-minerals:groups}{plagioclasa}.

\textbf{20.} La composición de la \omterm{def:b1:rock-forming-minerals:groups}{plagioclasa} (su contenido en anortita).

\textbf{21.} Un \omterm{def:b1:rock-forming-minerals:groups}{anfíbol} con un \omterm{def:b1:optical-mineralogy:extinction}{ángulo de extinción} de \ang{17}: hornblenda (\ang{15}--\ang{19}), no un \omterm{def:b1:rock-forming-minerals:groups}{piroxeno}
(\ang{38} para el diópsido).

\textbf{22.} $\qty{38000}{nm}\times0.020 = \qty{760}{nm}$: verde de segundo \omterm{def:b1:optical-mineralogy:interference-colour}{orden}.

\textbf{23.} Está cortado perpendicularmente a un \omterm{def:b1:optical-mineralogy:optic-axis}{eje óptico}: sus dos índices son iguales para la luz que viaja según el eje,
$\Delta = 0$, y permanece oscuro como un grano isótropo.

\textbf{24.} A cuarzo, B diópsido, C biotita, D \omterm{def:b1:rock-forming-minerals:groups}{plagioclasa}, E hornblenda: una roca ígnea intermedia con algo de cuarzo, como
una cuarzodiorita.

\textbf{25.} El espesor de la sección es \textbf{\qty[text-series-to-math=true]{38}{\micro m}}.
\end{solution}
